\section{A concrete nonabelian transfer: an \texorpdfstring{$A_5$}{A5} detector}
\label{sec:a5-transfer}

The earlier packets already propose $S_3$ and completely split-prime targets.
Here is a more demanding test case, together with a fully proved reduction to
one specific missing estimate.  The reduction uses the repository's one-sided
Landau mechanism and the now available RH for $\zeta$.  It does not assume
Artin's holomorphy conjecture.  The finite group calculations are independently
checked by the accompanying standard-library Python program
\texttt{a5\_detector.py}.  The arithmetic prime-sum estimate below is not proved.

Let $L/\Q$ be a Galois extension with group $A_5$.  Choose subgroups
$A_4,D_5,C_5$ of orders $12,10,5$, where $D_5$ denotes the dihedral group of
order $10$, and put
\[
 K_5=L^{A_4},\qquad K_6=L^{D_5},\qquad K_{12}=L^{C_5}.
\]
The subscripts are the degrees over $\Q$.  Write $\pi_j$ for the permutation
character of $A_5$ on the cosets defining $K_j$.

\begin{center}
\begin{tabular}{c|rrrrr}
Conjugacy class&$1$&$2A$&$3A$&$5A$&$5B$\\\hline
Class size&$1$&$15$&$20$&$12$&$12$\\
$\pi_5$&$5$&$1$&$2$&$0$&$0$\\
$\pi_6$&$6$&$2$&$0$&$1$&$1$\\
$\pi_{12}$&$12$&$0$&$0$&$2$&$2$\\\hline
$W=\pi_5+\pi_6+\pi_{12}-2$&$21$&$1$&$0$&$1$&$1$
\end{tabular}
\end{center}

\begin{lemma}[Two exact character identities]
\label{lem:a5-identities}
The function $W$ is nonnegative and has group average $1$.  The regular
character satisfies
\[
 \operatorname{reg}_{A_5}=4\pi_5+2\pi_6+3\pi_{12}-8.
\]
Consequently the following identity of meromorphic functions holds:
\begin{equation}
 \zeta_L(s)\zeta(s)^8
 =\zeta_{K_5}(s)^4\zeta_{K_6}(s)^2\zeta_{K_{12}}(s)^3.
 \label{eq:a5-zeta-identity}
\end{equation}
\end{lemma}

\begin{proof}
For a subgroup $H$, the coset character at $g$ is
$\#\{x\in A_5:x^{-1}gx\in H\}/|H|$.  This gives the table, and the
regular-character identity follows in each of its five columns.  The mean of
$W$ is $(21+15+12+12)/60=1$.  Artin formalism sends a permutation character on
$G/H$ to $\zeta_{L^H}$, sends sums to products, and sends the regular character
to $\zeta_L$.  Applying these identities first in $\Re s>1$, where Euler
products converge absolutely, and then continuing meromorphically proves
\eqref{eq:a5-zeta-identity}.  This use of Artin formalism does not assert the
holomorphy of individual nontrivial Artin factors; see, for example,
\cite{gun-hazra-sahu,ltz2024}.
\end{proof}

\subsection{A positive Euler coefficient at every prime power}

Define
\begin{equation}
 \mathcal F(s)=
 \frac{\zeta_{K_5}(s)\zeta_{K_6}(s)\zeta_{K_{12}}(s)}{\zeta(s)^2},
 \qquad
 D(s)=-\frac{\mathcal F'(s)}{\mathcal F(s)}.
 \label{eq:a5-detector-function}
\end{equation}
For a rational prime $p$, choose an inertia subgroup $I_p$ and a lift $g_p$
of its Frobenius element in the decomposition group.  Set
\begin{equation}
 a(p^k)=\log p\,\frac1{|I_p|}\sum_{h\in I_p}W(g_p^k h)
 \quad(k\geq1),
 \qquad a(n)=0\quad\hbox{otherwise}.
 \label{eq:a5-coefficients}
\end{equation}
The expression is independent of the choices.  Every summand in
\eqref{eq:a5-coefficients} is nonnegative, and
\[
 0\leq a(n)\leq21\Lambda(n),\qquad
 D(s)=\sum_{n\geq1}\frac{a(n)}{n^s}\quad(\Re s>1).
\]
Indeed, the projection onto inertia invariants is
$|I_p|^{-1}\sum_{h\in I_p}\rho(h)$, and its trace after applying
$\rho(g_p)^k$ is the average in \eqref{eq:a5-coefficients}.  The same identity
holds for a virtual character by linearity.  Expanding the local logarithmic
derivative gives the displayed Dirichlet series.  This is also the local
definition in \cite{gun-hazra-sahu}.  Thus neither ramified primes nor prime
powers are being discarded.  At an unramified prime the formula simplifies to
$a(p^k)=W(g_p^k)\log p$; using only $W(g_p)$ at all powers would be incorrect.

For example, at unramified primes the local factors of $\mathcal F$, with
$t=p^{-s}$, are
\[
\begin{array}{c|c}
1&(1-t)^{-21}\\
2A&(1-t)^{-1}(1-t^2)^{-10}\\
3A&(1-t^3)^{-7}\\
5A,5B&(1-t)^{-1}(1-t^5)^{-4}.
\end{array}
\]
These follow by computing the logarithmic derivative using the values of
$W(g^k)$, or directly from the three permutation actions.

\subsection{The precise missing estimate and its implication}

Each Dedekind zeta function has a simple pole at $1$, so $\mathcal F$ has a
simple pole there.  Let $b_{\mathcal F}$ be the constant term defined by
\[
 D(s)=\frac1{s-1}+b_{\mathcal F}+O(s-1),
 \qquad c_{\mathcal F}=b_{\mathcal F}-1,
\]
and define, for $x\geq1$,
\begin{equation}
 R_W(x)=\sum_{n\leq x}a(n)\left(\frac1n-\frac1x\right)
           -\log x-c_{\mathcal F}.
 \label{eq:a5-remainder}
\end{equation}
This specifies the constant without an unknown normalization.  In terms of
Euler--Kronecker constants it is
$c_{\mathcal F}=-1-\gamma_{K_5}-\gamma_{K_6}-\gamma_{K_{12}}+2\gamma$.

\begin{theorem}[A single nonabelian arithmetic target]
\label{thm:a5-criterion}
Assume RH for the Riemann zeta function.  Suppose that for every
$\varepsilon>0$ there are $C_\varepsilon,X_\varepsilon>0$ such that
\begin{equation}
 R_W(x)\leq C_\varepsilon x^{-1/2+\varepsilon}
 \qquad(x\geq X_\varepsilon).
 \tag{$A_5$-bound}\label{eq:a5-missing-bound}
\end{equation}
Then $\zeta_L$, $\zeta_{K_5}$, $\zeta_{K_6}$, and $\zeta_{K_{12}}$ satisfy
their Riemann hypotheses.  No holomorphy hypothesis for the individual
irreducible Artin $L$-functions is needed.
\end{theorem}

\begin{proof}
Fix $1/2<\beta<1$.  Choose $0<\varepsilon<\beta-1/2$; by increasing the
threshold, \eqref{eq:a5-missing-bound} implies
$R_W(x)\leq x^{\beta-1}$.  Put $E(v)=R_W(e^v)$ for $v\geq0$, and extend it
by zero for $v<0$.  The coefficient bound above
implies $E(v)=O(1+v^2)$ for $v\geq0$, so in particular $E$ is measurable and
has the growth allowed by \texttt{New.IsPolar}.  Direct integration in
$\Re s>1$ gives
\begin{align}
 \int_0^\infty E(v)e^{(1-s)v}\,dv
 &=\int_1^\infty R_W(x)x^{-s}\,dx\notag\\
 &=\frac{D(s)}{s(s-1)}-\frac1{(s-1)^2}
                         -\frac{c_{\mathcal F}}{s-1}.
 \label{eq:a5-transform}
\end{align}
To check the first term, integrate each ramp from $n$ to $\infty$:
\[
 \int_n^\infty\left(\frac1n-\frac1x\right)x^{-s}\,dx
 =\frac{n^{-s}}{s(s-1)}.
\]
At $s=1$ the Laurent expansion of the first term in
\eqref{eq:a5-transform} is
$(s-1)^{-2}+(b_{\mathcal F}-1)(s-1)^{-1}+O(1)$, so the two prescribed
subtractions cancel both poles.  Elsewhere in $\Re s>0$ the transform is
meromorphic with at most simple poles.  If $m(\rho)$ is the signed order of
$\mathcal F$ at $\rho\ne1$, its residue is
\begin{equation}
 r(\rho)=\frac{m(\rho)}{\rho(1-\rho)}.
 \label{eq:a5-residue}
\end{equation}
The sign here is essential: $-\mathcal F'/\mathcal F$ has residue $-m(\rho)$.

RH ensures that the denominator $\zeta(s)^2$ in
\eqref{eq:a5-detector-function} has no zero in $\Re s>1/2$.  Away from $1$,
all three numerator factors are holomorphic, whence $m(\rho)\geq0$ in that
half-plane.  For real $\beta<u<1$, \eqref{eq:a5-residue} is therefore
nonnegative.  At $u=1$ the corrected transform is holomorphic, and for real
$u>1$ the Euler products give $m(u)=0$.  The one-sided bound and these real
residue signs are exactly the input of
\texttt{New.IsPolar.landau} in \texttt{Lib/Analytic/Landau.lean}.
It follows that $r(\rho)=0$ for $\Re\rho>\beta$.  Thus $\mathcal F$ has no
zero there.  Since $\beta>1/2$ was arbitrary, this excludes all its zeros
with real part greater than $1/2$.

In this half-plane, away from $1$, the order of $\mathcal F$ is the sum of
the three nonnegative orders of its numerator zeta functions.  Their zeros
are therefore all excluded individually.  Equation
\eqref{eq:a5-zeta-identity} then excludes zeros of $\zeta_L$ as well.
The functional equations of the four Dedekind zeta functions reflect any
nontrivial zero with real part below $1/2$ to one above $1/2$.  This proves
the conclusion.
\end{proof}

The theorem is a sufficient criterion, not a proof of
\eqref{eq:a5-missing-bound}.  Its useful feature is the separation of tasks:
the analytic implication, treatment of real zeros, ramification, prime powers,
and denominator poles are complete.  A new short/enough/long-forms construction
must now reproduce the coefficients $W(\operatorname{Frob}_p^k)$ and prove
the displayed upper bound.  The source's current congruence and quadratic
weights do not supply it.

\subsection{A scoped optimality statement}

Consider only the normalized three-field family
\[
 W_{a,b,c}=a\pi_5+b\pi_6+c\pi_{12}+(1-a-b-c)\cdot1,
 \qquad a,b,c\geq1.
\]
The normalization gives mean $1$; the lower bounds require each of the three
Dedekind factors to be detected with coefficient at least $1$.  Its identity
value is
\[
 W_{a,b,c}(1)=1+4a+5b+11c\geq21.
\]
Thus any nonnegative member has peak at least $21$.  The table proves that
$a=b=c=1$ attains peak $21$.  This is an exact optimum in this declared family,
also for real coefficients.  It is not an assertion of global optimality
among all subfields, all kernels, or all analytic reductions.  The direct
regular-character detector has peak $60$ at the same mean $1$.

\subsection{Why fixed abelian congruence weights cannot replace this input}

\begin{theorem}[An exact mean-square obstruction]
\label{thm:a5-blindness}
Let $B/\Q$ be any finite abelian extension with group $H$, and let
$f:H\to\C$ be arbitrary.  For primes unramified in $LB$, write $g_p,h_p$
for the two Frobenius components.  Then
\[
 \lim_{X\to\infty}\frac1{\pi(X)}
 \sum_{\substack{p\leq X\\p\text{ unramified in }LB}}
 |W(g_p)-f(h_p)|^2
 =7+\frac1{|H|}\sum_{h\in H}|f(h)-1|^2\geq7.
\]
Equality holds for $f\equiv1$.  In particular, no fixed finite collection of
Dirichlet characters can represent the weight $W$ at all unramified primes.
\end{theorem}

\begin{proof}
$A_5$ has no nontrivial abelian quotient.  Therefore $L\cap B=\Q$ and
$\operatorname{Gal}(LB/\Q)=A_5\times H$.  Chebotarev turns the prime average
into the uniform group average.  The table gives
\[
 \frac1{60}\sum_{g\in A_5}W(g)=1,
 \qquad
 \frac1{60}\sum_{g\in A_5}(W(g)-1)^2
 =\frac{400+20}{60}=7.
\]
Expanding $|W-f|^2=|(W-1)+(1-f)|^2$, the mixed term vanishes after the
$A_5$ average; the displayed identity follows.  Any finite collection of
Dirichlet characters factors through one finite cyclotomic, hence abelian,
extension, so the final assertion follows as well.
\end{proof}

This obstruction concerns coefficient representations through a fixed finite
abelian extension.  It is not a logical independence theorem, does not rule
out a new argument using RH, and does not cover moduli changing with the
prime-sum scale.  Such changing moduli would in turn require new uniform
estimates, whereas the source's final bounds fix its modulus first.  The
general phenomenon that $A_5$ splitting behavior is not determined by
congruences is classical; the exact variance $7$ above is the calculation for
this particular detector.

\subsection{A completely specified field}

One can take $L$ to be the splitting field of
\[
 P(X)=X^5+20X+16.
\]
Here is a short Galois certificate, reproduced exactly by the companion
program.  The polynomial discriminant is
$1\,024\,000\,000=32\,000^2$, so the Galois group lies in $A_5$.
Modulo $3$, $P=X^5-X+1$ has no root.  Its remainders modulo the three monic
irreducible quadratics $X^2+1$, $X^2+X+2$, and $X^2+2X+2$ are respectively
$1$, $X+1$, and $X+1$.  A reducible quintic has a factor of degree at most
$2$, so $P$ is irreducible modulo $3$.  Modulo $7$,
\[
 P=(X+2)(X+3)(X^3+2X^2-2X-2),
\]
and the cubic has no root.  Thus the Galois group contains elements of orders
$5$ and $3$.  Its order divides $60$ and is divisible by $15$, leaving
$15,30,60$.  A group of order $15$ is cyclic by Sylow's theorems and cannot
embed in $S_5$; an index-two subgroup of $A_5$ would give a nontrivial map to
$C_2$, impossible since $A_5$ is generated by elements of order $3$.
Hence the group has order $60$ and is $A_5$.

This identifies a fixed, genuinely nonabelian zeta-function target.  No
claim is made that its RH has been proved by the present calculation.  Nor
does this report claim priority for the polynomial or Artin factorization.
The contribution is the explicit positive detector, its exact analytic
certificate and scoped optimization, and the quantified obstruction to
replacing its nonabelian coefficients by fixed congruence data.
